% digest.tex -- cumulative weekly scan for novel control-theory research.
% Updated automatically every Friday; newest week first. Bibliography: refs.bib.
\documentclass[11pt,a4paper]{article}
\usepackage[margin=2.5cm]{geometry}

\usepackage{acronym}
% The acronym environment typesets a list, so it lives in the document body
% (straight after \maketitle), not here.

\newcommand{\ts}[1]{{\textnormal{#1}}}
\newcommand{\diff}[2]{{\frac{d #1}{d #2}}}
\newcommand{\pdiff}[2]{{\frac{\partial #1}{\partial #2}}}
\newcommand{\ie}{\emph{i.e.},\ }
\newcommand{\eg}{\emph{e.g.}}
\newcommand{\cf}{\emph{c.f.}\ }
\newcommand{\Nset}{\mathbb{N}}
\newcommand{\Rset}{\mathbb{R}}
\newcommand{\Cset}{\mathbb{C}}
\newcommand{\mc}{\mathcal}
\newcommand{\md}{\mathds}
\newcommand{\mb}{\mathbf}
\newcommand{\mbb}{\mathbb}
\newcommand{\mf}{\mathfrak}
\newcommand{\ms}{\mathscr}
\newcommand\st{\colon~}
\newcommand{\ov}{\overline}

\usepackage{tikz}
\usepackage{pgfplots}
\pgfplotsset{every axis/.append style={line width=1pt}}
\pgfplotsset{every tick label/.append style={font=\tiny}}
\pgfplotsset{compat=newest}
\usetikzlibrary{shapes, fit, decorations.markings, matrix, plotmarks, positioning, spy, patterns, shadows, calc, backgrounds, arrows.meta, arrows}
\usetikzlibrary{decorations.pathreplacing}
\usetikzlibrary{circuits.ee.IEC}
\usetikzlibrary{fit}
\usetikzlibrary{plotmarks,pgfplots.colorbrewer}
\usetikzlibrary{decorations.pathreplacing}
\usetikzlibrary{circuits.ee.IEC}
\usepackage[americanvoltages,siunitx]{circuitikz}

% Loaded before the theorem definitions: \newtheorem* needs amsthm.
\usepackage{cite}
\usepackage{amsmath,amssymb,amsfonts,mathtools}
\usepackage{amsthm}

\newtheorem{theorem} {Theorem}
\newtheorem{example} {Example}
\newtheorem{definition} {Definition}
\newtheorem{assumption} {Assumption}
\newtheorem{lemma} {Lemma}
\newtheorem{coro} {Corollary}
\newtheorem{proposition}{Proposition}
 \newtheorem*{remark}{Remark}

% subequations
\newcounter{subeqn} \renewcommand{\thesubeqn}{\theequation\alph{subeqn}}%
\newcommand{\subeqn}{%
  \refstepcounter{subeqn}% Step subequation number
  \tag{\thesubeqn}% Label equation
}

\usepackage{algorithmic}
\usepackage{graphicx}
\usepackage{textcomp}
\usepackage{xcolor}
\usepackage[linesnumbered,ruled,vlined]{algorithm2e}
\usepackage[shortlabels]{enumitem}
\usepackage{bm}
\usepackage{stfloats}
\usepackage{dsfont}

\usepackage[colorlinks=true,allcolors=blue!50!black]{hyperref} % clickable links to the papers

\title{Control Theory: Weekly Scan for Novel Work}
\author{Automated weekly scan}
\date{Last updated: none yet} % LAST-UPDATED (each weekly run rewrites this line)

\begin{document}
\maketitle

\section*{Acronyms}
% Alphabetical. Each weekly run adds any new acronym it uses, never a duplicate.
\begin{acronym}[MPC]
  \acro{LMI}{linear matrix inequality}
  \acro{MPC}{model predictive control}
  \acro{RL}{reinforcement learning}
\end{acronym}

\section*{Selection bar}
A paper is picked only if it clearly does at least one of the following:
introduces a new problem formulation, framework or system-theoretic concept;
brings a mathematical tool into control in a way that changes what can be proved or computed;
proves a fundamental limit, an impossibility result or a tight characterisation, or settles a known open problem;
removes a standing assumption, or changes scalability or the kind of guarantee qualitatively.
Excluded: extensions that reuse the same proof technique, less conservative \ac{LMI} conditions,
known methods (\eg, \ac{MPC} or \ac{RL}) applied to a new plant, and simulation-only claims.
arXiv preprints are assessed from their full text; journal articles from their abstracts only.

Sources: arXiv (eess.SY and math.OC, including cross-lists) and new articles in IEEE TAC, Automatica, SICON,
IEEE TCNS, MCSS, ESAIM:COCV and SIADS, plus control-related articles in SIOPT, Mathematical Programming and
Mathematics of Operations Research.

% ---------------------------------------------------------------------------
% Weekly entries, newest first. Each run inserts its section directly below
% the marker line and leaves the marker in place.
% ---------------------------------------------------------------------------
%%% NEW-WEEK-INSERTION-POINT %%%

\nocite{*}
\bibliographystyle{IEEEtran}
\bibliography{refs}
\end{document}
